\documentclass[11pt]{article}
\usepackage[margin=1in]{geometry}
\usepackage{amsmath,amssymb,amsthm}
\usepackage[hidelinks]{hyperref}
\emergencystretch=3em

\newcommand{\ca}{\operatorname{ca}}
\newcommand{\diam}{\operatorname{diam}}
\newcommand{\clu}{\operatorname{clust}}
\newcommand{\dist}{\operatorname{dist}}

\title{Literature-implied status for Oikhberg--Spinu's Schur-lattice remark}
\author{}
\date{2026-07-18}

\begin{document}
\maketitle

\paragraph{Status.}
This is a \emph{literature-implied answer}, not a new proof.  In
Oikhberg--Spinu, \emph{Domination of operators in the non-commutative
setting}, arXiv:1209.0699, Remark~\texttt{rem:not l\_1} in the Positive Schur
Property subsection says that an order-continuous atomic Banach lattice with
the Schur property need not be isomorphic to \(\ell_1\): if
\((E_n)\) are finite-dimensional lattices, then
\[
  E=\left(\sum_{n=1}^\infty E_n\right)_{\ell_1}
\]
has the Schur property, and the choice \(E_n=\ell_2^n\) is not isomorphic to
\(\ell_1\).  The remark then says that the authors do not know any Banach
lattice with the Schur property which is not isomorphic to an \(\ell_1\)-sum
of finite-dimensional spaces.

\paragraph{The supporting example.}
Kacena--Kalenda--Spurny, \emph{Quantitative Dunford-Pettis property},
arXiv:1110.1243, Example~10.1, constructs a Banach space \(X\) such that
\(X^*\) is a separable \(L\)-embedded space with the Schur property.  In the
proof, for each \(\alpha>0\) they define \(X_\alpha\) on \(c_0\) and compute
\[
  \|x^*\|_\alpha^*=\alpha\|x^*\|_1+\|x^*\|_\infty
  \qquad (x^*\in \ell_1).
\]
Then
\[
  X=\left(\bigoplus_{n\in\mathbb N} X_{1/n}\right)_{c_0},
  \qquad
  X^*=\left(\bigoplus_{n\in\mathbb N} X_{1/n}^*\right)_{\ell_1},
\]
and the paper states that this \(X^*\) has the Schur property.

Each \(X_\alpha^*\) is a Banach lattice under the usual coordinate order on
\(\ell_1\): if \(|u|\le |v|\), then both \(\|u\|_1\le \|v\|_1\) and
\(\|u\|_\infty\le \|v\|_\infty\).  The norm is equivalent to the usual
\(\ell_1\)-norm, so the space is complete; it is atomic and order-continuous
by the usual monotone-convergence argument on \(\ell_1\).  Therefore
\[
  E:=X^*=\left(\bigoplus_{n\in\mathbb N} X_{1/n}^*\right)_{\ell_1}
\]
is an atomic order-continuous Banach lattice with the Schur property.

\paragraph{Why it is not an \(\ell_1\)-sum of finite-dimensional spaces.}
Kalenda, \emph{Quantitative Schur property and measures of weak
non-compactness}, arXiv:2505.12893, defines the \(c\)-Schur property by
\[
  \ca(x_n)\le c\,\delta(x_n)
\]
for every bounded sequence \((x_n)\), where \(\delta(x_n)\) is the usual
measure of weak non-Cauchyness, equivalently the diameter of the set of
weak-star cluster points of \((x_n)\) in the bidual.  The same paper notes, in
Example~(2) after Corollary~\texttt{cor:qschur}, that if \(X\) is the space
from Kacena--Kalenda--Spurny Example~10.1, then \(X^*\) is \(L\)-embedded and
has the Schur property, but \(\operatorname{wk}(\cdot)\) and
\(\omega(\cdot)\) are not equivalent in \(X^*\); hence \(X^*\) fails the
quantitative Schur property.

On the other hand, Kalenda's Theorem~\texttt{T:ell1suma} says that, for
\(c\ge 1\), the \(c\)-Schur property is preserved by arbitrary \(\ell_1\)-sums.
Every finite-dimensional Banach space has the \(1\)-Schur property: bounded
sequences have compact norm closure, and the diameter of the norm cluster set
coincides with \(\ca(x_n)\), while the weak and norm cluster sets coincide in
finite dimension.  Thus every \(\ell_1\)-sum of finite-dimensional spaces has
the quantitative Schur property.

Finally, quantitative Schur is invariant under Banach-space isomorphism up to
changing the constant.  If \(T:Y\to Z\) is an isomorphism and \(Y\) has the
\(c\)-Schur property, then for \(z_n=Ty_n\),
\[
  \ca(z_n)\le \|T\|\ca(y_n)
  \le c\|T\|\delta_Y(y_n)
  \le c\|T\|\,\|T^{-1}\|\,\delta_Z(z_n).
\]
Consequently any space isomorphic to an \(\ell_1\)-sum of finite-dimensional
spaces must have the quantitative Schur property.  Since the above
Kacena--Kalenda--Spurny \(X^*\) fails quantitative Schur, it is not
isomorphic, even as a Banach space, to an \(\ell_1\)-sum of finite-dimensional
spaces.

\paragraph{Conclusion and scope.}
The space
\[
  E=\left(\bigoplus_{n\in\mathbb N}
  (\ell_1,\ \|x\|_{1/n}=(1/n)\|x\|_1+\|x\|_\infty)\right)_{\ell_1}
\]
is a Banach lattice with the Schur property, but it is not isomorphic to an
\(\ell_1\)-sum of finite-dimensional spaces.  This supplies a negative answer
to the ``we do not know of any'' remark in arXiv:1209.0699.  The supporting
papers do not explicitly cite that remark, so the status is
\texttt{literature\_implied\_answer} rather than
\texttt{literature\_already\_answered}.  No claim is made here about the main
non-commutative domination questions in the source paper.

\begin{thebibliography}{9}
\bibitem{OS}
T. Oikhberg and E. Spinu,
\emph{Domination of operators in the non-commutative setting},
arXiv:1209.0699.

\bibitem{KKS}
M. Kacena, O. F. K. Kalenda, and J. Spurny,
\emph{Quantitative Dunford-Pettis property},
arXiv:1110.1243.

\bibitem{Kalenda}
O. F. K. Kalenda,
\emph{Quantitative Schur property and measures of weak non-compactness},
arXiv:2505.12893.
\end{thebibliography}

\end{document}
